\section{Discussion}
\label{sec:discussion}

\paragraph{What the comparison was measuring.} The ordering in \tabref{tab:main} is
reproducible and, as a description of trained models under the usual recipe, correct.
What it does not license is the inference that explicit codes compute worse. On copy
the opposite holds: under a forced output length rotary reaches $0.996$ token
accuracy at $L{=}10$ and $0.984$ at $L{=}12$ against $0.987$ and $0.877$ for \nope{},
and with EOS removed from training rotary is perfect at $L{=}16$ while \nope{} falls
to $0.43$. The standard protocol was measuring where models stop.

\paragraph{Why the boundary lands on $L_{\mathrm{tr}}$.} A model trained at one
length sees exactly $L_{\mathrm{tr}}$ answer tokens in every example. If it can
address output positions, the rule ``emit EOS at index $L_{\mathrm{tr}}{+}1$'' is a
single parameter that fits the training distribution perfectly and is simply
undefined outside it; nothing in the objective penalises that. If it cannot address
positions, the same rule is not representable, and the model must fall back on
content or on a distance cue, both of which are fuzzier and later---which is what
the \nope{} and ALiBi profiles in \figref{fig:profile} show. The asymmetry is not
about the informativeness of the code but about whether a stop can be encoded as a
constant.

This is the mechanism behind the ``length attractor'' of \citet{newman2020eos}:
they showed that models trained to predict EOS stratify hidden states by position
and stall there, and we locate the consequence at the level of a cross-code
comparison. It is also why a mixture over lengths does not help the absolute code
(\secref{sec:moves}): the constant becomes a lookup with one entry per length, but
it is still a lookup, and the entry for the test length does not exist.

\paragraph{Relation to the no-PE recommendation.} \citet{kazemnejad2023impact}
compare five encodings on ten tasks and find \nope{} best. Our \tabref{tab:main}
reproduces that ordering on our tasks and our model size. The result here does not
overturn their measurement; it reinterprets what it measures. On tasks where the
answer has the same length as the input, an exact-match comparison is partly a
comparison of how crisply each code can place a stop. The practical corollary is
cheap: report forced-length accuracy alongside exact match, or evaluate with EOS
masked, and the comparison becomes about computation again.

\paragraph{Why ALiBi is the exception in the standard table.} ALiBi is the only
standard code that both leaves the residual stream alone and is unaddressable: its
bias orders output positions by distance from the query without naming any of them.
It therefore stops late and softly rather than at a constant, which is what its
$0.92$ exact match at $L{=}10$ reflects. We would predict that any code with a
monotone and unaddressable positional signal behaves similarly, and that ALiBi's
advantage would shrink under no-EOS training---a prediction we did not test.

\paragraph{What this does not show.} We study copying with a 3.2M-parameter model,
two seeds, and one training length per condition. \emph{sort}, \emph{reverse} and
\emph{prefix-XOR} are in \tabref{tab:main} but not in the decomposition, and the
permutation tasks may not decompose the same way, since a truncated answer there is
wrong at every position rather than at the tail. We also did not test whether the
effect survives scale, pretraining, or natural language, where answers rarely have
the same length as inputs and the boundary is cued by content rather than by a
count. The claim we are prepared to defend is the narrow one: on copy-like
same-length tasks, at this scale, the cross-code ranking on exact match is dominated
by termination, and it reverses when termination is removed.
